\documentclass[sigplan, nonacm = true, natbib = false]{acmart}

\usepackage{listings}

\lstset{
  basicstyle=\ttfamily\footnotesize,
  columns=fullflexible,
  keepspaces=true
}

\lstdefinestyle{metl}{
  language=[Objective]Caml,
  morekeywords={eff, handle, end, return, mask, forall, Effect, alias},
  morecomment=[l]{--},
  literate=
    {->}{{$\to$}}2
    {=>}{{$\Rightarrow$}}2
    {()}{{()}}2
}

\lstdefinestyle{koka}{
  morekeywords={fun,rec,handler,forall},
  sensitive=true,
  morecomment=[l]{--},
  morestring=[b]",
  literate=
    {->}{{$\to$}}2
    {()}{{()}}2
}

\usepackage[backend=bibtex, style=numeric]{biblatex}
\addbibresource{abstract.bib}
\renewcommand*{\bibfont}{\footnotesize}

%% Ghosts

\usepackage{halloweenmath}
\makeatletter
\newcommand*\forallghostraise{.16ex} % vertical nudge upwards
\newcommand*\forallscale{0.4}       % scale factor for ∀

\newcommand{\forallghost}{\mathord{\mathpalette\forallghost@{}}}
\newcommand{\forallghost@}[2]{%
  \ooalign{%
    % ghost outline
    \hfil$\m@th#1\mathpalette\raiseghost\relax$\hfil\cr
    % smaller forall symbol, raised
    \hfil\raise\forallghostraise\hbox{%
      \scalebox{\forallscale}{$\m@th#1\forall$}}%
    \hfil\cr
  }%
}

\makeatother


\newcommand{\myghost}{\mathpalette\raiseghost\relax}
\newcommand{\raiseghost}[2]{\raisebox{.5\depth}{\scalebox{0.9}{$#1\mathghost$}}}

\newcommand{\Hole}{{\myghost}}
\newcommand{\TopHole}{\forallghost} % implicit forall

%%

%% Keywords

\newcommand{\keyw}[1]{{\textbf{{\textsf{#1}}}}}
\newcommand{\Do}{\keyw{do}\;}
%%

\title{Practical programming with Modal Effect Types (Research Presentation)}
\author{}
\keywords{effect handlers, effect types, modal types, type inference,
  first-class polymorphism}

\begin{document}
\sloppy
\begin{abstract}
  Effect handlers are gaining in popularity and existing languages start
  supporting them. However, existing type-and-effect systems lack either
  backward compatibility or robustness.

  We overview our implementation of a toy ML-like language supporting effects
  trying to solve both of these shortcomings. To do so, it combines modal effect
  types and a new type inference algorithm for first-class polymorphism.
\end{abstract}
\maketitle
\section{Introduction}
\subsection{Existing type-and-effect systems}
To take effects into account in the type system, traditional languages such as
\textsf{Koka}~\cite{koka2014} annotate arrows with the effects the function can
perform. Thus, each function inherits, in its signature, the effects of its
callees. This becomes cumbersome when typing higher-order functions and is an
obstacle to backward compatibility. For instance, in \textsf{Koka}, the function
\textsf{map} has the following type, even though it does not perform effects
itself.
\begin{lstlisting}[style=koka]
  forall <a, b, e> . (list<a>, (a) -> e b) -> e list<b>
\end{lstlisting}
Syntactic sugar can simplify such signatures, as demonstrated by
Frank~\cite{frank2020}. But, relying on syntax makes the system fragile.

The language \textsf{Effekt}~\cite{effekt2020} answers this problem with a
\emph{capability} based system. However, to guarantee effect-safety, functions
become second-class objects, hindering backward compatibilty.

\subsection{Modal Effect Types}
\textsf{\textsc{Met}}~\cite{tang2025modal, tang2026modal} draws inspiration from
Multimodal Type Theory~\cite{mtt2020} and represents effects a computation can
perform with \emph{modalities}. For instance, the \textsf{map} function has the
following type, where the modality \textsf{[]} indicates that the function does
not perform any effect.
\begin{lstlisting}[style=metl]
  [](forall a b . (a -> b) -> list a -> list b)
\end{lstlisting}

As modalities require explicit boxing and unboxing, one needs an inference
algorithm to turn \textsf{\textsc{Met}} into a programming language. The
original paper~\cite{tang2025modal} provides such an algorithm. But it only
deals with modalities and requires explicit type application.

\subsection{Frost}
As noted in previous work on \textsf{\textsc{Met}}, elimination and introduction
of modalities remind instantiation and generalisation in type
inference~\cite[Section 7]{tang2025modal}. \textsf{Frost} is a new
bidirectionnal type inference algorithm for first-class polymorphism which makes
formal this parallel. This is possible because \textsf{Frost} relies on
\emph{skeleton} inference. Skeletons are types haunted by \emph{ghosts}
($\Hole$) and \emph{universal ghosts} ($\TopHole$). The algorithm will fill the
former with other skeletons, and the latter with quantifiers and modalities.

\section{Type inference for Modal Effect Types}
\subsection{Programming with \textsf{\textsc{Met}}}
We give an overview of \textsf{\textsc{Met}} with scoped
rows~\cite{scopedrows2005} as an effect structure (i.e.\ our modalities will
have scoped rows of effects). TO CONTINUE

All typing judgements take a row of effects as input called the \emph{ambient
effect context} which indicates the effects the current computation can
perform. For instance, the expression $\Do \textsf{yield } 1$ is well-typed in
the effect context $\textsf{gen int}$, where \lstinline[style=metl]{eff gen [a] = yield : a => unit}.
The brackets around the type parameter \textsf{a} indicate that it has kind
\textsf{Abs}: the kind of types whose meaning does not depend on the ambient
effect context. To avoid leakage, the types of inputs and outputs of operation
have kind \textsf{Abs}.

\emph{Modalities} $\mu$ can change the ambient effect contexts: either by
replacing it by an independant one (absolute modalities: $[E]$), or by removing
and adding some effects (relative modalities: $\langle L|D\rangle$). The
$\keyw{mod}_\mu$ and $\keyw{let mod}_\mu$ constructs respectively introduce and
eliminate modalities. For example, this expression performs explicit
subeffecting of a variable $x:[](\textsf{int} \to\textsf{int})$:
\[
\left(\keyw{let mod}_{[]} y = x \keyw{ in } \keyw{mod}_{[E]}y\right) : [E](\textsf{int} \to\textsf{int}).
\]

Handlers also modify the effect context of the expression they encompass, to add
the effects they handle. In an empty effect context, this expression is
well-typed:
\[
\keyw{handle } (\Do \textsf{yield } 1) \keyw{ with }\{\textsf{yield } k\; n
\Rightarrow n \} : \textsf{int}.
\]
\subsection{Parallel between modalities and quantifiers}
For subtyping, \textsf{Frost} relies on the broom
\subsection{Examples}
\subsection{Cooperative concurrency}
Cooperative concurrency effect is an effect with two operations with signatures
that look like \lstinline[style=metl]{fork : (unit -> unit) => unit} and
\lstinline[style=metl]{suspend : unit => unit}. Implementing a scheduler for
such an effect serves as a test for compositionality~\cite[Section~6]{frank2020}
and expressivity~\cite[Section~2.10]{tang2025modal} of an effect system.

To implement this example in \textsf{\textsc{Met}}, one can add effect
polymorphism to the language. This allows to give pure types to the argument of
fork.
\begin{lstlisting}[style=metl]
  eff co (e : Effect) = fork : [co e, e]{unit} => unit
                      | suspend : unit => unit
\end{lstlisting}

However, if the scheduler uses another effect \lstinline[style=metl]{queue} to
manage the running functions, it has type

\begin{lstlisting}[style=metl]
  forall (e : Effect) .
    [qproc e] ([co e, qproc e]{unit} -> unit),
\end{lstlisting}
in presence of the following definitions 
\begin{lstlisting}[style=metl]
  type proc (e : Effect) = Proc of [queue (proc e), e]{unit}
  alias qproc e = queue (proc e), e.
\end{lstlisting}
This type signature reminds traditional type-and-effect systems and looks like a
desugared version of a \textsf{Frank} type. Moreover, the implementation relies
on a \emph{modality-parametrised handler}~\cite[Section~3.5]{tang2026modal} to
give a type with an absolute modality to the continuation, making it even more
cumbersome.

REWRITE These problems motivate the search of other ways to express higher-order
effects. We have implemented a promising approach: relax kind restrictions on
operations' operands, but restrict effect subtyping in presence of operations
with non-pure operands. This means that $l, l',D\leq l', l,D$ if, and only if
$l'$ is not higher-order, and that one can perform an operation from a
higher-order effect if, and only if it is the first entry in the effect
context. We conjecture that these restrictions guarantee type safety by ensuring
that values coming from higher-order effects stay in a fixed effect
context. With these rules, the scheduler has the following type, where
\lstinline[style=metl]{queue} and \lstinline[style=metl]{co} are higher-order,
and it does not rely on modality-parametrised handler.
\begin{lstlisting}[style=metl]
  [queue {unit}](<co>{unit} -> unit)
\end{lstlisting}

For now, we lack a proof of safety and an expressivity result relating this
system with effect polymorphism.
\subsection{Some limitations}
Our current unification rules have a syntactical treatment of quantifiers. Thus
we can unify neither \lstinline[style=metl]{[l]<l'>(unit -> unit)} and
\lstinline[style=metl]{[l', l](unit -> unit)},
nor \lstinline[style=metl]{forall a b . a * b}
and \lstinline[style=metl]{forall b a . a * b}, while they are both subtype of eachother in both cases.

Lack of kind inference
\begin{lstlisting}[style=metl]
let find p l =
  handle iter (fun x -> if mask<gen> (p x) then do yield x else ()) l with
  | return _ => None
  | yield x k => Some (x)
  end
\end{lstlisting}
\subsection{Future work}
\printbibliography
\end{document}


